% ===========================================================================
\chapter{The fewest steps: breadth-first search}\label{ch:bfs}

We have a graph. Now: what is the best way from one node to another? This
chapter and the next two answer that question three times, each time a
little better. All three algorithms share one idea. Starting from the
source, they explore the graph outwards, keeping a list of nodes they know
how to reach but haven't explored yet, called the \emph{frontier}. They
differ only in which frontier node they explore next.

\section{What a search finds}

Every search gives back the same kind of answer: the path it found, as a
list of nodes from the source to the target, and the total cost of that
path. It also gives back two things that describe its work, so that we can
compare the algorithms and draw them: the nodes it explored, in the order
it explored them, and how it reached each node.

That last part is the trick that makes all these searches cheap. When the
search first reaches node $n$ from node $m$, it writes down
``$n$ came from $m$'': $m$ is $n$'s \emph{parent}. It never needs to store
whole paths. When it reaches the target, it follows the parent links back
to the source, and that is the path, backwards.

\begin{newfile}{search.nex}
-- search.nex: finding a way from one node to another. Breadth-first
-- search, Dijkstra's algorithm and A*.

intern Road_Graph

-- What a search found, and how much work it did to find it.
class Search_Result
create
  make(name: String, nodes: Array[Integer], total: Real,
    order: Array[Integer], tree: Map[Integer, Integer]) do
    algorithm := name
    path := nodes
    cost := total
    explored := order
    parents := tree
  end
feature
  algorithm: String
  -- The nodes of the way found, from start to finish; empty if there is
  -- no way.
  path: Array[Integer]
  -- What the search was minimising, added up along the path.
  cost: Real
  -- Every node the search finished with, in the order it did so.
  explored: Array[Integer]
  -- How the search reached each node it saw: parents.get(n) is the node
  -- it came to n from. Together these links form a tree.
  parents: Map[Integer, Integer]

  found(): Boolean do
    result := not path.is_empty
  end
end

-- The way from `source` to `target`, following the parent links back
-- from `target` and then turning the list around.
function trace_back(parents: Map[Integer, Integer], source, target: Integer): Array[Integer] do
  let backwards: Array[Integer] := [target]
  let at: Integer := target
  from
  until
    at = source
  do
    at := parents.get(at)
    backwards.add(at)
  end
  result := backwards.reverse
ensure
  starts_at_source: result.get(0) = source
  ends_at_target: result.get(result.length - 1) = target
end
\end{newfile}

Two more helpers, which we'll use in contracts and checks: whether a list
of nodes really is a route through the graph (every step must be an
edge), and the total cost of a route, for any way of costing an edge.

\begin{append}{search.nex}
-- Whether `path` is a way through the graph: every step is an edge.
function is_route(g: Road_Graph, path: Array[Integer]): Boolean do
  result := not path.is_empty
  from
    let i: Integer := 1
  until
    not result or i >= path.length
  do
    result := g.edge_between(path.get(i - 1), path.get(i)) /= nil
    i := i + 1
  end
end

-- The total of `cost` over the steps of a route.
function route_cost(g: Road_Graph, path: Array[Integer], cost: Function(Edge): Real): Real
require
  is_a_route: is_route(g, path)
do
  from
    let i: Integer := 1
  until
    i >= path.length
  do
    if ?g.edge_between(path.get(i - 1), path.get(i)) as e then
      result := result + cost(e)
    end
    i := i + 1
  end
end
\end{append}

\code{route_cost}'s last parameter is a \emph{function}:
\code{Function(Edge): Real} is the type of a function that takes an edge
and gives back a number. Passing it in lets one routine add up lengths,
times, or anything else. Here are the two ways of costing an edge we will
use, each as a function that returns the costing function:

\begin{append}{search.nex}
-- Two ways to measure an edge: by its length, or by how long it takes.
function by_distance(): Function(Edge): Real do
  result := fn (e: Edge): Real do result := e.length end
end

function by_time(): Function(Edge): Real do
  result := fn (e: Edge): Real do result := e.seconds end
end
\end{append}

\code{fn (e: Edge): Real do ... end} is an \emph{anonymous function}: a
function written as a value, without a name. Inside it, \code{result} is
the anonymous function's own result.

\section{Breadth-first search}

The simplest search explores the frontier in the order it found it: first
in, first out, like a queue at a shop. Start with just the source in the
queue. Repeatedly take the node at the front of the queue; if it is the
target, stop; otherwise put each of its neighbours that hasn't been seen
before at the back of the queue, with the node as its parent.

\begin{center}
\begin{tikzpicture}[
  n/.style={circle, draw=accent, fill=white, inner sep=0pt, minimum size=5.5mm, font=\scriptsize},
  e/.style={color=cm}]
  \node[n, fill=st!30] (s) at (0,0) {0};
  \foreach \i/\x/\y in {a/1.4/0.9, b/1.4/-0.9, c/-1.4/0} \node[n, fill=accent!25] (\i) at (\x,\y) {1};
  \foreach \i/\x/\y in {d/2.8/1.5, f/2.8/0.2, g/2.8/-1.4, h/-2.8/0.8, k/-2.8/-0.8} \node[n, fill=accent!10] (\i) at (\x,\y) {2};
  \node[n] (m) at (4.2,0.9) {3};
  \draw[e] (s) -- (a) (s) -- (b) (s) -- (c) (a) -- (d) (a) -- (f) (b) -- (g) (c) -- (h) (c) -- (k) (d) -- (m) (f) -- (m);
  \node[font=\small, text=cm, align=left, anchor=west] at (5,0) {each node's number:\\how many steps it is\\from the source};
\end{tikzpicture}
\end{center}

Nodes come out of the queue in rings: the source, then every node one step
from it, then every node two steps away, and so on. A node is put in the
queue the first time it is seen, from a node one ring further in. So when
the target comes out, the parent links lead back to the source along a
path with \emph{the fewest steps} there is. That is what breadth-first
search finds.

Removing the first element of a Nex array shifts all the others down, so
instead of removing anything, we keep an index, \code{front}, of the next
node to take. The nodes before it have all been explored.

\begin{append}{search.nex}
-- Breadth-first search: the way with the fewest steps.
--
-- The search keeps a queue of nodes to visit, first in, first out. It
-- starts with just the source. It takes the node at the front of the
-- queue and puts every neighbour it has not seen before at the back, and
-- repeats until it takes out the target. Nodes come out in order of how
-- many steps they are from the source: all the nodes 1 step away, then
-- all those 2 steps away, and so on. So the first time the search reaches
-- the target, it has reached it in as few steps as possible.
function breadth_first(g: Road_Graph, source, target: Integer): Search_Result
require
  source_on_roads: g.has_node(source)
  target_on_roads: g.has_node(target)
do
  let queue: Array[Integer] := [source]
  let front: Integer := 0
  let seen: Set[Integer] := #{source}
  let parents: Map[Integer, Integer] := {}
  let explored: Array[Integer] := []
  let found: Boolean := false
  from
  until
    found or front = queue.length
  do
    let node: Integer := queue.get(front)
    front := front + 1
    explored.add(node)
    if node = target then
      found := true
    else
      across g.edges_from(node) as e do
        if not seen.contains(e.target) then
          seen.add(e.target)
          parents.put(e.target, node)
          queue.add(e.target)
        end
      end
    end
  end
  let path: Array[Integer] := []
  if found then
    path := trace_back(parents, source, target)
  end
  result := create Search_Result.make("breadth-first", path,
    (path.length - 1).max(0).to_real, explored, parents)
ensure
  valid_path: not result.found() or is_route(g, result.path)
end
\end{append}

If the queue runs out before the target turns up, the target can't be
reached, and the path is empty. The cost of a breadth-first path is its
number of steps.

The postcondition says that whatever the search finds is a real route:
``either nothing was found, or the path is a route''. It is checked every
time the search runs.

\section{Watching a search}

To see what the searches do, we'll draw them: the map, then a line from
every explored node to its parent in purple, which shows the part of the
graph the search covered, and the path it found in blue. This try-it file
holds the drawing code, and the next few chapters' try-its use it with
\code{intern Try_Pictures}.

\begin{tryit}{try_pictures.nex}
-- try_pictures.nex: drawing a search on the map, for the try-its of the
-- next few chapters.

intern Map_View
intern Search
intern Helpers

-- Draw the map, then the part of the graph that `found` explored in
-- purple, then the way it found in blue, and save the picture.
function picture_of_search(view: Map_View, found: Search_Result, file: String) do
  let island: Island_Map := view.island
  view.draw_map()
  view.set_pen(create Color.make_rgba(150, 60, 200, 170), 4.0)
  across found.explored as n do
    if found.parents.contains_key(n) then
      let a: Point := island.point(found.parents.get(n))
      let b: Point := island.point(n)
      view.canvas.line(a.x, a.y, b.x, b.y)
    end
  end
  view.line(island.points_of(found.path), create Color.make(40, 110, 220), 5.0)
  view.save(file)
end

-- A one-line summary of a search, for driving.
function summary(g: Road_Graph, found: Search_Result): String do
  result := found.algorithm + ": " + found.explored.length.to_string + " nodes explored, "
    + (found.path.length - 1).to_string + " steps, "
    + distance_text(route_cost(g, found.path, by_distance())) + ", "
    + duration_text(route_cost(g, found.path, by_time()))
end
\end{tryit}

Now a search from the beach car park, on the south-west coast, to the
lighthouse, on the northern cape. A place isn't a node of the graph, so we
snap each one to its nearest node first.

\begin{tryit}{try_bfs.nex}
intern Try_Pictures

let island: Island_Map := create Island_Map.load(user_path("kestrel_island.map"))
let g: Road_Graph := create Road_Graph.build(island, Mode.Driving)
if ?island.find_place("beach") as a and ?island.find_place("lighthouse") as b then
  let source: Integer := g.nearest_node(island.point(a.node_id))
  let target: Integer := g.nearest_node(island.point(b.node_id))
  let found: Search_Result := breadth_first(g, source, target)
  console.print_line(summary(g, found))
  picture_of_search(create Map_View.make(island, 1000, 700), found, "bfs.png")
end
\end{tryit}
\tryoutput{try_bfs}
\mappicture{bfs}

The purple spreads out evenly in every direction from the car park, along
every road, until it touches the lighthouse. The path it found goes all the
way round the west of the island on the Coast Road.

Is that the best way? It is the way with the fewest steps, but a step is
not a fixed distance. Look back at the map file: the nodes along the Coast
Road are about 320 metres apart, while the lanes have one every 170 metres
or so, and the town's streets one every 100. So a route along the Coast
Road has few, long steps, and breadth-first search likes it for that
reason alone. The road up the middle of the island, by Lake Road and Hill
Road, is shorter in metres, but it has more steps.

Breadth-first search is the right tool when every step costs the same,
like moves in a puzzle or links between web pages. Roads are not like
that. We need a search that counts metres, or seconds.
